% =====================================================================
\section{HTTPS и OTA: защищённое устройство}
% =====================================================================
\lessonintro
  {проверять подлинность сервера и обновлять прошивку без USB, не оставляя устройство «кирпичом»}
  {урок 5 (Wi-Fi), урок 11 (Flash и NVS)}
  {устройство безопасно получает данные и может обновляться по сети}
  {доступ к защищённой Wi-Fi-сети; для OTA --- стабильное питание}

\subsection{TLS не равен «порт 443»}

HTTPS защищает содержимое канала и проверяет, что устройство разговаривает с тем сервером. Для
проверки ESP32-S3 нужна корневая цепочка сертификата (Root CA) и правильное время: сертификат имеет
срок действия. Если просто вызвать \code{client.setInsecure()}, соединение заработает, но устройство
перестанет отличать ваш сервер от подмены --- это допустимо только для краткого лабораторного опыта.

\begin{figure}[H]
\centering
\begin{tikzpicture}[y=-0.8cm,font=\sffamily\scriptsize]
  \node[blk,fill=blkcpu,minimum width=3cm] (e) at (0,0) {ESP32-S3};
  \node[blk,fill=blkrf,minimum width=3cm] (s) at (9,0) {Сервер :443};
  \draw[thick,gray] (e.south) -- (0,7.3);
  \draw[thick,gray] (s.south) -- (9,7.3);
  \draw[arr,accent] (0,1.3) -- node[above]{0. NTP: узнать текущее время (до TLS)} (7,1.3);
  \draw[arr,accent] (0,2.4) -- node[above]{1. ClientHello: поддерживаемые шифры} (9,2.4);
  \draw[arr,okgreen] (9,3.5) -- node[above]{2. ServerHello + цепочка сертификатов} (0,3.5);
  \fill[esp!30] (-0.15,3.7) rectangle (0.15,4.9);
  \node[anchor=east,align=right] at (-0.3,4.3) {3. проверка: подпись Root CA,\\ имя сервера, срок действия};
  \draw[darr] (0,5.6) -- node[above]{4. обмен ключами --- общий секрет} (9,5.6);
  \draw[darr,esp] (0,6.8) -- node[above]{5. зашифрованный HTTP: GET / \ldots} (9,6.8);
\end{tikzpicture}
\caption{Упрощённая схема TLS-соединения. Без точного времени шаг 3 не проверит срок действия
сертификата, а без Root CA --- подлинность сервера}
\end{figure}

\codecap{Синхронизация времени и проверка Root CA}
\codeinput{l13_https}

Сертификат берите у центра сертификации сервера и обновляйте с запасом до истечения. Не кладите в
прошивку пароль Wi-Fi, API-токен и приватный ключ публичного проекта. Для устройства, которое будет
продаваться, включают Secure Boot, Flash Encryption и подписанные образы в ESP-IDF.

\subsection{OTA: два слота вместо одной ставки}

OTA записывает новый образ в неиспользуемый раздел Flash. После перезапуска загрузчик запускает новый
слот. Если в загрузчике включён откат (rollback), слот становится постоянным только после того, как
новая прошивка подтвердила работоспособность. Это требует
таблицы разделов с двумя OTA-слотами и разделом \code{otadata}. Нельзя перезаписывать текущую работающую
прошивку.

\codecap{Минимальный OTA-сервис Arduino для локальной сети}
\codeinput{l13_ota}

Этот вариант удобен для лаборатории, но не является удалённым промышленным обновлением: добавьте
пароль, проверку подписи/хэша образа, HTTPS-скачивание и журнал версии. После прошивки устройство
должно самостоятельно проверить датчики, сеть и ключевые задачи, а затем подтвердить загрузчику
новую версию. Если прошивка падает до подтверждения, загрузчик возвращается к прошлому слоту --- но
только при включённой опции \code{CONFIG\_BOOTLOADER\_APP\_ROLLBACK\_ENABLE}. В стандартной сборке
Arduino-ESP32 она выключена, поэтому для настоящего изделия откат настраивают в ESP-IDF.

\subsection{Разбор ошибки: HTTPS или OTA не запускается}

При ошибке TLS сначала выведите текущую дату: почти всегда причина --- устройство не дождалось NTP,
не видит DNS или имеет старый Root CA. Не заменяйте диагностику \code{setInsecure()}. При обрыве OTA
не выключайте питание вручную: двухслотовая разметка сохраняет старую прошивку. Если устройство
больше не отвечает после эксперимента, зайдите в загрузчик кнопками BOOT и RST из урока 0 и прошейте
USB-кабелем.

\begin{summary}
\begin{itemize}[leftmargin=*]
  \item TLS требует Root CA и корректных часов; \code{setInsecure()} не решает задачу безопасности.
  \item Безопасное OTA --- это два слота, проверка образа и подтверждение успешного запуска.
  \item Обновление не должно останавливать контроль нагрузки и не должно зависеть от бесконечного ожидания сети.
\end{itemize}
\end{summary}

\begin{selfcheck}
\begin{enumerate}
  \item Зачем устройству точное время для HTTPS?
  \item Почему \code{setInsecure()} не годится для готового устройства?
  \item Зачем OTA-обновлению два слота прошивки?
\end{enumerate}
\end{selfcheck}

\begin{tasks}
\begin{enumerate}
\item \lv{1} Объясните своими словами, чем HTTPS отличается от HTTP и что проверяет сертификат.
\item \lv{1} Выведите время UTC после NTP: без него HTTPS не должен запускаться.
\item \lv{2} Добавьте в веб-страницу из урока 5 номер версии, дату сборки и свободное место Flash.
\item \lv{2} Выполните запрос HTTPS с проверкой сертификата и проверьте, что при неверном времени он завершается ошибкой.
\item \lv{3} Составьте checklist OTA: питание, версия, хэш, тест после перезапуска, откат и журнал.
\item \lv{4} Сделайте обновление по сети: плата скачивает прошивку, проверяет и подтверждает её после перезагрузки, а при сбое откатывается.
\end{enumerate}
\end{tasks}
